\documentclass[../../main.tex]{subfiles}
\begin{document}

\section{黏贴 2 维胞腔}\label{section:ch11-黏贴2维胞腔}

我们已经用两种方式算出了环面 $T=S^1\times S^1$ 的基本群：其一是通过考虑标准覆叠映射 $p\times p:\mathbb{R}\times\mathbb{R}\to S^1\times S^1$ 并使用提升工具。其二是应用乘积空间基本群的基本定理。本节再给出一个计算环面的基本群的办法。

若将覆叠映射 $p\times p$ 限制在单位矩形上，我们得到一个商映射 $\pi:I^2\to T$。它将 $\operatorname{Bd}I^2$ 映到子空间 $A=(S^1\times b_0)\cup(b_0\times S^1)$ 上，并且将 $I^2$ 中的其他部分一一地映射到 $T-A$ 上，其中 $A$ 为两个圆周的束。因此 $T$ 可以被想象为将正方形 $I^2$ 的边界黏在 $A$ 上。

将平面上某多边形区域的边黏到另外一个空间上是构造新空间的常用方法。在此，我们给出如何计算这种空间的基本群的办法。这个办法有许多漂亮的应用。

\begin{theorem}\label{theorem:ch11-attaching-2-cell}
设 $X$ 是一个 Hausdorff 空间，$A$ 为 $X$ 中的一个道路连通的闭子空间。假设存在一个连续映射 $h:B^2\to X$ 将 $\operatorname{Int}B^2$ 一一地映射到 $X-A$ 上，并且将 $S^1=\operatorname{Bd}B^2$ 映入 $A$ 中。设 $p\in S^1$ 并且 $a=h(p)$，$k:(S^1,p)\to(A,a)$ 为限制 $h$ 而得到的映射。则内射诱导的同态
\begin{align*}
i_*:\pi_1(A,a)\to\pi_1(X,a)
\end{align*}
是一个满射，并且它的核是 $\pi_1(A,a)$ 中包含 $k_*:\pi_1(S^1,p)\to\pi_1(A,a)$ 的像的最小正规子群。
\end{theorem}

有时我们说，$X$ 的基本群是从 $A$ 的基本群中“消除”类 $k_*[f]$ 而得到的，其中 $[f]$ 生成 $\pi_1(S^1,p)$。

\par
\begin{figure}[htbp]
\centering
\includegraphics[width=0.6\textwidth]{image/fig-72-1.png}
\caption{}\label{figure:ch11-fig-72-1}
\end{figure}

\par
\begin{figure}[htbp]
\centering
\includegraphics[width=0.4\textwidth]{image/fig-72-2.png}
\caption{}\label{figure:ch11-fig-72-2}
\end{figure}

\begin{proof}
\textbf{第一步。} 原点 $\mathbf{0}$ 为 $B^2$ 的圆心，设 $x_0$ 为 $X$ 中的点 $h(\mathbf{0})$。如果 $U$ 是 $X$ 中的开集 $U=X-x_0$，我们证明 $A$ 为 $U$ 的一个形变收缩核。见图 72.1。


设 $C=h(B^2)$，令 $\pi:B^2\to C$ 是限制 $h$ 的值域所得到的映射。考虑映射
\begin{align*}
\pi\times\operatorname{id}:B^2\times I\to C\times I,
\end{align*}
这是一个闭映射，因为 $B^2\times I$ 是紧的并且 $C\times I$ 是 Hausdorff 空间。因此，它是一个商映射。它的限制
\begin{align*}
\pi':(B^2-\mathbf{0})\times I\to(C-x_0)\times I
\end{align*}
也是一个商映射，这是由于它的定义域是 $B^2\times I$ 中的开集，并且是相对于 $\pi\times\operatorname{id}$ 的饱和开集。有一个从 $B^2-\mathbf{0}$ 到 $S^1$ 的形变收缩，经由商映射 $\pi'$ 它诱导了一个从 $C-x_0$ 到 $\pi(S^1)$ 的形变收缩。我们将这个形变收缩扩充到 $U\times I$ 上，并且在形变的过程中要求 $A$ 中的每一个点都保持不动。因此 $A$ 是 $U$ 的一个形变收缩核。

由此推出，$A$ 到 $U$ 的内射诱导出基本群之间的一个同构。从而，我们的定理归结为验证以下论断：

设 $f$ 是一条回路，它所在的类生成 $\pi_1(S^1,p)$。则由 $U$ 到 $X$ 的内射诱导一个满同态
\begin{align*}
\pi_1(U,a)\to\pi_1(X,a),
\end{align*}
其核是包含回路 $g=h\circ f$ 所在的类的最小正规子群。

\textbf{第二步。} 为方便起见，我们先考虑由内射诱导的同态 $\pi_1(U,b)\to\pi_1(X,b)$，其中基点 $b$ 不在 $A$ 中。

设 $b$ 为 $U-A$ 中的任何一点。将 $X$ 写成两个开集 $U$ 与 $V=X-A=\pi(\operatorname{Int}B^2)$ 之并。这时 $U$ 是道路连通的，因为它以 $A$ 为形变收缩核。由于 $\pi$ 为商映射，它在 $\operatorname{Int}B^2$ 上的限制也是一个商映射，从而是一个同胚。因此，$V$ 是单连通的。集合 $U\cap V=V-x_0$ 同胚于 $\operatorname{Int}B^2-\mathbf{0}$，因此它是道路连通的并且它的基本群为无限循环群。由于 $b$ 为 $U\cap V$ 的一个点，推论 70.4 蕴涵着由内射诱导的同态
\begin{align*}
\pi_1(U,b)\to\pi_1(X,b)
\end{align*}
是一个满射，并且它的核是包含无限循环群 $\pi_1(U\cap V,b)$ 的像的最小正规子群。

\textbf{第三步。} 我们再将基点改回点 $a$，并证明定理。

设 $q$ 为 $B^2$ 的一个点，它是 $\mathbf{0}$ 到 $p$ 的线段的中点，取 $b=h(q)$。这时 $b$ 便是 $U\cap V$ 的一个点。设 $f_0$ 是以 $q$ 为基点的 $\operatorname{Int}B^2-\mathbf{0}$ 中的一条回路，要求它表示这个空间的基本群的一个生成元，则 $g_0=h\circ f_0$ 是 $U\cap V$ 中的一个以 $b$ 为基点的回路，并且表示 $U\cap V$ 的基本群的一个生成元。见图 72.2。


由第二步得知，由内射诱导的同态 $\pi_1(U,b)\to\pi_1(X,b)$ 是一个满射，并且它的核是包含回路 $g_0=h\circ f_0$ 所在的类的最小正规子群。为了对以 $a$ 为基点的回路得到相应的结论，我们论证如下：

设 $\gamma$ 是 $B^2$ 中从 $q$ 到 $p$ 的一条直线道路，$\delta=h\circ\gamma$ 是 $U$ 中从 $b$ 到 $a$ 的一条道路。由道路 $\delta$ 分别在 $U$ 和 $X$ 中诱导的两个同构（都记作 $\hat{\delta}$）与由内射诱导的同态在以下图表中可交换：
\begin{align*}
\begin{matrix}
\pi_1(U,b)&\longrightarrow&\pi_1(X,b)\\
\hat{\delta}\Bigg\downarrow&&\Bigg\downarrow\hat{\delta}\\
\pi_1(U,a)&\longrightarrow&\pi_1(X,a)
\end{matrix}
\end{align*}
因此，由内射诱导的从 $\pi_1(U,a)$ 到 $\pi_1(X,a)$ 的同态是满的，并且它的核是包含元素 $\hat{\delta}([g_0])$ 的最小正规子群。

回路 $f_0$ 表示 $\operatorname{Int}B^2-\mathbf{0}$ 的以 $q$ 为基点的基本群的一个生成元。从而回路 $\bar{\gamma}*(f_0*\gamma)$ 表示 $B^2-\mathbf{0}$ 的以 $p$ 为基点的基本群的一个生成元。因此，它道路同伦于 $f$ 或者 $f$ 的逆。不妨设前者成立。将此道路同伦再复合映射 $h$，我们有 $\bar{\delta}*(g_0*\delta)$ 在 $U$ 中与 $g$ 道路同伦。于是 $\hat{\delta}([g_0])=[g]$，从而定理成立。
\end{proof}

在定理的证明中，对于 $B^2$ 本身的性质，我们并没有用到许多。事实上，如果我们将 $B^2$ 换成一个与之同胚的空间 $B$，并且以 $\operatorname{Bd}B$ 的同胚像代替 $S^1$，则上述定理的相应结论依然成立。这样的空间 $B$ 称为 2-维胞腔（2-cell）。上述定理中的空间 $X$ 可以想象为将一个“2-维胞腔”黏贴到 $A$ 上。后面我们将对这种情形进行更为规范的讨论。

\subsection*{习题}

\begin{exercise}\label{exercise:ch11-72-01}
设 $X$ 为 Hausdorff 空间，$A$ 是一个道路连通的闭子空间。假设连续映射 $h:B^n\to X$ 将 $S^{n-1}$ 映入 $A$ 中，并且将 $\operatorname{Int}B^n$ 一一地映为 $X-A$。设 $a$ 为 $h(S^{n-1})$ 的一个点。若 $n>2$，试说明由内射诱导的从 $\pi_1(A,a)$ 到 $\pi_1(X,a)$ 的同态有哪些性质。
\end{exercise}

\begin{exercise}\label{exercise:ch11-72-02}
设 $X$ 是由一个正规的道路连通空间 $A$ 与单位球 $B^2$ 的无交并并且借助于连续映射 $f:S^1\to A$ 给出的贴附空间。（见第 35 节习题 8。）证明 $X$ 满足定理 72.1 中的假设条件。何处用到了 $A$ 的正规性？
\end{exercise}

\begin{exercise}\label{exercise:ch11-72-03}
设 $G$ 是一个群，$x$ 是 $G$ 中的一个元素，$N$ 为 $G$ 的包含 $x$ 的最小正规子群。证明：若存在一个正规的道路连通空间使得其基本群同构于 $G$，则存在一个正规的道路连通空间，其基本群同构于 $G/N$。
\end{exercise}

\end{document}